\chapter{第十七族和第十八族：抢电子与不爱反应}

\begin{kaichang}
今天晚上回家，请你找三样东西：厨房里的一包食盐、洗手台上的一支含氟牙膏、药箱里的一小瓶碘伏。它们一个管调味，一个管牙齿，一个管伤口，看上去八竿子打不着。

可是化学家会说：这三样东西“沾亲带故”。盐里有氯，牙膏里有氟，碘伏里有碘——这三种元素在周期表上住在同一列，倒数第二列，第 17 族。同一家族的成员，一个你每天放心吃，一个你主动往嘴里送，一个你往破皮的伤口上抹。而这个家族的“本来面目”却相当吓人：氯气在第一次世界大战里当过化学武器，氟气能跟几乎任何东西起火，纯碘的蒸气能把皮肤烧出泡。

为什么同一族的元素，换一身“装束”，就从危险品变成了日用品？它们右边的邻居、周期表最后一列第 18 族，故事更离奇：那一整列元素曾经被所有教科书宣布“永远不会和任何物质反应”，这个信念稳稳当当撑了将近七十年——然后被一个下午的计算和一烧瓶气体打破了。

先猜一猜：让“惰性气体绝不反应”这条铁律破产的，是一场意外的实验事故，还是一次有预谋的推理？
\end{kaichang}

\section{先看这一家的颜色和脾气}

第 17 族有个正式名字叫\textbf{卤素}，意思是“成盐的元素”——因为它们和金属凑在一起，产物多半是盐。这一族排在最前面的是四位：氟、氯、溴、碘。先看它们单质的模样，这是任何一本手册都查得到的硬事实：

\begin{itemize}[itemsep=2pt]
\item \textbf{氟气}（\chem{F_2}）：淡黄色气体，脾气暴烈到连盛放它都麻烦。
\item \textbf{氯气}（\chem{Cl_2}）：黄绿色气体，有刺鼻的消毒水味——游泳池边那股味道就和它有关。
\item \textbf{溴}（\chem{Br_2}）：深红棕色液体，是常温常压下唯一呈液态的非金属单质，一打开瓶盖就挥发出呛人的红棕色蒸气。
\item \textbf{碘}（\chem{I_2}）：紫黑色固体，稍微加热就变成紫色蒸气，凝华回来又结成亮晶晶的小片。
\end{itemize}

注意到没有？从氟到碘，单质的状态从气体、液体一路变到固体。分子越来越大，分子之间的拉扯越来越强，所以要更高的温度才能让它们“散伙”（第 22 章会专门讲分子之间的这种吸引力，这里先记下现象）。

脾气方面，四位是同一个路数——都极爱抢电子——但力气差别明显。氟最凶：它几乎能和所有元素直接反应，连黄金都不放过；也正因为太凶，人类直到 1886 年才由法国化学家莫瓦桑第一次用电解的办法单独制出氟气。氯稍收敛，但仍然能和铁、铜、氢气直接反应。溴温和一些。碘最不积极，和铁粉反应都得加热才肯动手。

有一个干净利落的办法能直接比出力气大小：让两种卤素“抢同一个东西”。把氯水滴进无色的溴化钠溶液里，溶液很快变成橙黄色——那是被“顶”出来的溴。氯把溴离子手里的电子抢走了：

\[\chem{Cl_2} + \chem{2NaBr} \rightarrow \chem{2NaCl} + \chem{Br_2}\]

反过来，把溴水滴进氯化钠溶液，什么也不会发生。力气大的抢力气小的，从不倒过来。用同样的方法可以验证：溴能把碘顶出来，而碘谁也顶不动。氟、氯、溴、碘，抢电子的力气依次递减——这不是背下来的口诀，是一瓶一瓶试出来的事实。

再看最右边那一列。第 18 族是氦、氖、氩、氪、氙、氡，如今叫\textbf{稀有气体}（老名字叫“惰性气体”，这个名字后来被证明起错了，证据层再讲）。它们的单质全是无色无味的气体，而且出了名的“什么都不干”：不成键、不反应，单原子独来独往。你早就见过它们工作的样子——霓虹灯。灯管里充不同的稀有气体，通电后发出不同颜色的光：氖是标志性的橘红色，氩是淡紫色，氙气灯发出接近日光的刺眼白光，汽车大灯和电影放映机都爱用它。发光的原理你在第 9 章见过：电流把电子踢上高楼层，电子落回低楼层时把能量以特定颜色的光吐出来。

\section{差一个就住满的楼层}

为什么这一族都爱抢电子，而隔壁那一族什么都不爱干？镜头推进到原子的最外层。

回忆第 9 章的电子楼层和第 10 章的周期表地图：元素的化学性格几乎完全由最外层电子数决定。卤素原子的最外层都住着 7 个电子，离住满只差 1 个；稀有气体原子的最外层恰好住满（氦是第一层住满 2 个，其余都是最外层 8 个）。

先说卤素。当一个氯原子遇到一份“多出来”的电子——比如钠原子那个急着脱手的最外层电子——把这份电子收进自己第三层的最后一个空位，整个体系的能量会降低。第 10 章已经强调过：原子没有愿望，不会“为了凑八个电子”而努力；真实的原因是能量。收进这份电子要放出能量，放出来的能量如果能盖过整个过程中其他要付的能量账，变化就会自己发生。对卤素来说，这笔账几乎总是划算的，所以它们逮着机会就抢。

抢到电子之后，氯原子就不再是氯原子了——它变成\textbf{氯离子}，写作 \chem{Cl^-}。这个负离子的外层住满了 8 个电子，能量低、脾气安稳，再也不到处抢东西。这就是开场那个谜团的钥匙：\textbf{氯气和食盐里的氯，是同一种元素的两种电子状态，化学性格天差地别}。第 10 章用一句话点过这件事，这一章我们要把它挖到底。

那稀有气体为什么不掺和？把同样的能量账算一遍就明白了。想\textbf{给}氩原子塞一个电子？它的第三层已经住满，新电子只能住进更高的第四层——能量不降反升，没人愿意做亏本买卖。想\textbf{从}氩原子身上夺走一个电子？它对现有电子抓得很紧，夺走第一个电子要付的能量（电离能）在所有元素里排前列，抢它的电子太贵。两头都亏，于是两千年……不，几十亿年下来，它们几乎什么反应都不参加。

注意这里的措辞：是“不划算”，不是“被禁止”。自然规律里没有“禁止稀有气体反应”这一条，只有一本能量账。既然只是账本问题，那么只要找到一笔能算得过账的生意，稀有气体也可能被拉下水。记住这句话，证据层要用。

\begin{qiaoliang}
抢电子的力气，能不能用数字说话？可以。科学家把“一个气态原子收进一个电子时放出的能量”叫作\textbf{电子亲和能}，单位是每摩尔千焦（\ud{}{kJ/mol}）。数字越大，说明这个原子收电子放出的能量越多、抓得越紧。查表可得：氯约 $348$、溴约 $325$、碘约 $295$（单位均为 \ud{}{kJ/mol}）——一路递减，正好和“氯能顶出溴、溴能顶出碘”的实验顺序吻合。

为什么越往下力气越小？想想第 11 章讲过的那场拉锯战：新来的电子是被原子核吸引的，但它入住的楼层一层比一层高，离核越来越远，中间隔着层层内层电子像棉被一样屏蔽核的吸引力。碘的新电子要住进第五层，核的“手”伸到那么远已经没什么力气了，放出的能量自然少。反过来，同一个道理也解释了稀有气体：它们外层住满，亲和能约等于零甚至为负——塞电子进去不但不放能，还要倒贴。
\end{qiaoliang}

\section{抢走电子之后，世界换了模样}

现在把符号请出来，写一笔完整的账。钠和氯气相遇：

\[\chem{2Na} + \chem{Cl_2} \rightarrow \chem{2NaCl}\]

每个钠原子交出自己的外层电子，每个氯原子（氯气分子里成对存在，写成 \chem{Cl_2}）收进一个电子。产物是钠离子 \chem{Na^+} 和氯离子 \chem{Cl^-}，它们靠正负电荷的吸引整整齐齐堆成食盐晶体——这堆法本身的能量账（离子键和晶格能）留到第 18 章细讲，本章只需要记住结果：\textbf{反应之后，两种危险品都消失了}。遇水就炸的钠没有了，灼伤肺部的氯气没有了，剩下的是可以撒在菜里的白色晶体。

这就是化学里最值得早一天明白的道理之一：\textbf{元素的性质，不等于它的单质的性质，也不等于它在某种化合物里的表现}。身份证（质子数）没变，电子状态和邻居变了，性格就全变了。氯离子不但无害，还是你身体的必需品：胃液里的盐酸靠它配平电荷，体液的渗透压靠它和钠离子一起维持，你出汗流失的盐里一半就是它。

反过来，氯气为什么是危险品？正因为分子里那两个氯原子还没“吃饱”——它们共用最外层电子（共价键，第 19 章细讲），但抓住更多电子的能量账仍然划算。氯气进入呼吸道，就在你的湿润组织里继续它那套抢电子、和水反应的化学，生成的东西破坏细胞。同一个元素，吃饱了的安稳，没吃饱的伤人，就这么简单。

再来看看这个家族的三份“正经工作”，你就能体会人类怎样把卤素的脾气变成工具：

\textbf{消毒剂}。自来水厂里通入少量氯气（或加入次氯酸盐），它在水中生成的次氯酸分子能钻进细菌体内，氧化破坏维持生命的关键分子，把微生物一锅端。家里常用的 84 消毒液，有效成分是次氯酸钠 \chem{NaClO}，是同一类化学的稀释版。这里的关键词是\textbf{剂量}：杀灭微生物所需的浓度，和对人体明显有害的浓度之间，隔着很大的余量，水厂的任务就是把余氯控制在“够用又安全”的窄窗口里——你闻到的淡淡氯味，正是这个窗口还守着的信号。

\textbf{含氟牙膏}。牙齿表面的珐琅质主要是一种叫羟基磷灰石的矿物（大致成分 \chem{Ca_{10}(PO_4)_{6}(OH)_2}），口腔细菌产酸时会把它一点点溶解，这就是蛀牙的开始。牙膏里的氟离子会替换掉矿物表面的羟基，生成氟磷灰石——同样是钙盐晶体，但耐酸能力明显更强，酸要溶解它得费更大的劲。等于给牙齿表面换了一层更抗腐蚀的“装甲”。当然，氟也是剂量见高下的典型：适量防蛀，长期过量反而会让牙齿出现斑块，所以儿童牙膏的氟含量和用量都有讲究。

\textbf{卤素灯}。老式白炽灯的钨丝在高温下慢慢蒸发变细，灯越烧越暗、丝越烧越脆。往灯泡里充入少量碘或溴（配上一半是氩之类的保护气），事情就变了：蒸发出来的钨原子飘到温度较低的泡壳附近，和碘结合成碘化钨气体；碘化钨飘回灼热的灯丝附近又被高温拆开，钨原子重新回到灯丝上，碘脱身再去跑下一趟。这叫卤钨循环——碘当了免费搬运工，钨丝寿命延长，灯可以在更高温度下工作，所以汽车大灯和舞台灯又小又亮。你在第 31 章会知道，这其实是一场被温度差牵着走的可逆反应。

\begin{shiyan}
\textbf{让碘现形：一粒米饭的显色实验}

材料：一小瓶碘伏（药店有售，就是擦伤口用的那种棕色消毒液）、一小块馒头或米饭、一片生土豆、一根棉签。

步骤：
\begin{enumerate}[itemsep=1pt]
\item 用棉签蘸一点碘伏，分别涂在一小粒米饭和土豆切面上。
\item 等十几秒，观察颜色变化。
\item 再用棉签蘸清水涂在另一粒米饭上作对照。
\end{enumerate}

观察点：涂过碘伏的地方会从棕黄色变成蓝黑色甚至接近黑色，而清水对照没有变化。这个蓝黑色是碘分子钻进淀粉分子的螺旋“隧道”里形成的——淀粉是米饭和土豆的主要成分（第 46 章还会和它打交道）。这个显色反应灵敏到能检出极微量的碘，化学家和生物学家用了一百多年。

安全提示：碘伏是家用消毒剂，少量沾到皮肤无妨，但\textbf{不要口服、不要弄进眼睛}；实验后的食物直接扔掉，不要再吃。
\end{shiyan}

\section{一个下午打破的七十年神话}

\begin{zhengju}
从 1894 年氩被发现起，化学家试过各种办法想让稀有气体参与反应，全部失败。到 20 世纪中叶，“稀有气体的外层电子结构决定了它们绝对不反应”已经写进所有教科书，几乎没人再试——反正试了也白试。

1962 年，在加拿大不列颠哥伦比亚大学工作的英国化学家尼尔·巴特利特（Neil Bartlett）正在研究一种脾气极坏的红色固体：六氟化铂 \chem{PtF_6}。这是已知最强的“电子强盗”之一。一次实验中他发现，\chem{PtF_6} 竟然把氧气分子 \chem{O_2} 的电子抢走了，生成了一种此前没人见过的盐。

换作别人，记录了这个奇怪结果也许就放过去了。巴特利特多做了一步：他去查表。从氧气分子身上夺走一个电子需要约 \ud{1175}{kJ/mol} 的能量。他记得另一个数字——从氙原子身上夺走一个电子需要约 \ud{1170}{kJ/mol}。两个数字几乎一样。

于是他推理：既然 \chem{PtF_6} 抢得动氧分子的电子，而氙原子的电子并不比氧分子的难抢，那么 \chem{PtF_6} 就应该抢得动氙的电子。外层住满？那本账在 \chem{PtF_6} 这样的强盗面前未必算不过来。这是一次\textbf{有预谋}的尝试，依据是两条公开数据和一个类比。他把无色的氙气和深红色的 \chem{PtF_6} 蒸气通进同一个容器——当天就得到了一种橙黄色固体。第一个稀有气体化合物，六氟合铂酸氙，就这样诞生了。事后分析还带来一个小修正：这种固体的真实成分比最初写的式子复杂（可以近似看作 \chem{XeF^+} 与 \chem{PtF_5^-} 的组合），最初的化学式写简单了——这个结果同样被同行如实修正，科学的自我纠偏照常运转。

消息传开，全世界实验室蜂拥而上。几个月内，二氟化氙 \chem{XeF_2}、四氟化氙 \chem{XeF_4} 相继合成；不久之后氪也被拉下了水。七十年的“绝对”就此作废，“惰性气体”这个名不副实的称呼也逐步让位给“稀有气体”。

这个故事值得记住的方法只有一条：巴特利特没有把“稀有气体不反应”当成自然界的法律，而是当成一句关于能量账的经验总结。经验总结的意思就是——只要账本变了，结论就可以变。他做的事不是推翻权威，而是把两句话并排放在桌上：“氧分子电离要 1175”“氙原子电离要 1170”，然后问了一个任何人都能问的问题。区别只在于，他去查了表。
\end{zhengju}

\section{一勺盐里，有多少个吃饱了的氯}

感受一下“抢电子”这件事在你厨房里发生的规模。一平勺食盐大约 5 克。食盐（\chem{NaCl}）的摩尔质量约 \ud{58.5}{g/mol}，所以这一勺约有 $5/58.5 \approx 0.085$ 摩尔的氯化钠。每摩尔包含 \ud{6.02\times 10^{23}}{} 个结构单元（第 5 章的阿伏伽德罗常数），算下来：

\[0.085 \times 6.02\times 10^{23} \approx 5\times 10^{22}\]

这一勺盐里大约有 $5\times 10^{22}$ 个氯离子——每个都是某个氯原子抢到电子后安稳度日的样子。这个数字有多大？全世界人口约 80 亿，即 $8\times 10^{9}$。$5\times 10^{22}$ 是它的约六万亿倍：就算给地球上每个人分一万亿个氯离子，你这一勺盐还发不完。而每一个氯离子背后，都发生过一次放能的电子转移。化学变化不是稀有事件，它在你家的盐罐子里以天文数字的规模静静完成着。

\section{这套说法什么时候会失灵}

本章的模型很好用，但任何模型都有边界，这一章尤其要把边界画清楚，因为这一族的历史本身就是一部“边界史”。

第一，\textbf{“外层住满就万事不理”是趋势，不是定律}。氙、氪的外层同样住满，照样被六氟化铂和氟气攻陷。今天已知的稀有气体化合物有几十种，但都集中在氙和氪，而且“对手”几乎总是氟、氧这类最凶的电子强盗；氦、氖、氩至今没有在常规条件下形成真正的化合物（2016 年，科学家在理论预言后用超高压让氦和钠形成了一种奇特结构 \chem{Na_2He}，但那需要比地心还高的压强）。所以准确的说法是：稀有气体反应性极低、自上而下依次“松动”，而不是绝对为零。

第二，\textbf{“抢电子力气氟氯溴碘递减”只在比较同一件事时成立}。换成别的账本就可能有例外。最著名的例外藏在电子亲和能里：氟的电子亲和能（约 \ud{328}{kJ/mol}）反而比氯（约 \ud{348}{kJ/mol}）略低——氟原子太小，新电子挤进狭小的第二层，电子之间的排斥吃掉了一部分收益。可氟气仍是实际反应中最凶的电子强盗，因为它在别处把账找了回来（章末挑战层里细说）。结论是：趋势箭头很好用，但每个箭头背后都是几笔账的总和，换一笔账，箭头就可能拐弯。

第三，\textbf{单质的危险程度推不出元素在化合物里的表现}。这条本章已经反复用，但值得单列为边界，因为它是生活中被违反得最多的一条：“某种元素的单质有毒”推不出“含这种元素的任何物质都有毒”，反过来，“某种元素是生命必需”也推不出“它的单质可以吃”。判断一种物质，永远要看它此刻的电子状态和邻居，而不是查元素名。

\begin{wuqu}
“都是消毒的，混在一起用效果加倍。”——这是真实出过人命的误解。84 消毒液的有效成分是碱性的次氯酸钠 \chem{NaClO}，洁厕灵一类清洁剂则含较强的酸。两者一旦混合，会发生放出氯气的反应：次氯酸根在酸性条件下和氯离子反应，把 \chem{Cl_2} 重新释放出来。卫生间空间小、通风差，新闻里屡有混合清洁剂导致氯气中毒送医的事故。这里不需要你记任何配方和比例，只需要记住原则：\textbf{化学品管不管“合得来”，不看货架上的分类，看粒子层面的账}。两种东西都贴着“清洁”标签，不代表它们的分子愿意和平共处。家用清洁剂一次只用一种，用完冲净再换下一种。
\end{wuqu}

\section{回到洗手台边的三样东西}

现在可以给开场的三样东西做个了断。

食盐里的氯：抢到电子、外层住满的氯离子，能量低、脾气安稳，和钠离子整整齐齐码成晶体。它不再抢任何东西，所以你能每天吃它——你的身体还离不开它。

含氟牙膏里的氟：同样是抢到电子后的氟离子，只是它的工作岗位在你的牙齿表面，把珐琅质的表层换成更耐酸的氟磷灰石。它不再腐蚀任何东西——抢电子的力气，已经在出厂前的化学反应里用完了。

碘伏里的碘：这是还没抢到电子的单质碘，配在溶液里。它的“没吃饱”正是它的用处——碘能破坏微生物的关键分子，所以能杀菌；也正因为它还有脾气，所以纯碘晶体和浓碘酊会刺激皮肤，得稀释了用，用完的饭粒不能吃。

同一族的元素，抢电子前是危险品，抢电子后是日用品，抢了一半、控制得当，就是工具。这一族的全部故事，就写在这一笔电子账里。

\section{抢与让之外，还有更大的世界}

第 12 章带你看了急着交电子的第 1、2 族，这一章看了抢着收电子的第 17 族和两不相干的第 18 族。一个交、一个收、一个不理——周期表右半边的故事似乎已经讲完了？

还早。周期表中间偏右，住着四位既不把电子完全交出去、也不能只靠抢解决问题的元素：氧、硫、氮、磷。氧抢电子的力气仅次于氟，却主要用“共享”的方式解决问题；氮的三个原子靠共享电子抱成一团，稳到闪电和细菌才能拆开，而拆开它们养活了全人类的化肥工业；磷排起了生命能量账本的队。这四位是\textbf{生命骨架元素}——你身体干重的绝大部分由它们和碳、氢组成。第 14 章，我们就去拜访这个撑起生命和环境的“四边形”。

至于卤素留下的悬念——氯分子里两个氯原子共用电子是怎么回事？氯化钠晶体里离子靠什么堆在一起？那是化学键的故事，第 17、18 章见。

\begin{tiaozhan}
挑战层（跳过不影响主线）：既然氟的电子亲和能比氯还低，凭什么说氟是最凶的卤素？答案是：反应算的是总账，电子亲和能只是其中一笔。以卤素抢电子的常见场景为例，总账至少有三笔：①拆开卤素分子要付的能量——\chem{F_2} 分子里的键意外地弱（拆开氟分子比拆开氯分子便宜得多，可能和两个小原子挤在一起、未共用电子对互相排斥有关）；②原子收电子放出的能量（电子亲和能）——这一笔氟确实略输给氯；③生成的负离子被水或晶体“接住”时放出的能量——\chem{F^-} 个子小、电荷集中，和水分子或晶格结合得特别牢，放出特别多。三笔加总，氟遥遥领先。所以“氟最强”和“氟的电子亲和能不是最高”并不矛盾——它们说的是不同粒度的账。这也是理解一切周期趋势的正确姿势：先问“比较的是哪本账”。
\end{tiaozhan}

\section*{练一练}

\begin{sikaoti}
\item 画两幅“楼层图”：左边画氯原子（17 个电子），右边画氯离子（18 个电子），标出最外层住满与没住满的差别，并用一句话说明：为什么左边的会出现在化学武器里，右边的会出现在你的菜汤里。图旁注明：楼层图是帮助思考的表示法，不是电子的真实照片。
\item 把氯水滴进无色的碘化钠溶液，预测你会看到什么颜色变化？用“电子搬家”的语言（谁把电子给了谁）描述发生了什么，并解释为什么反过来做（碘水滴进氯化钠溶液）不会有现象。
\item 巴特利特的关键一步是发现“氙的电离能约等于氧分子的电离能”。假设将来发现某种惰性程度未知的物质，它的电离能和氪几乎相等，你能据此做出什么预言？设计一个能检验这个预言的实验思路（只说思路，不必动手）。
\item 一勺 5 克的食盐约有 $5\times 10^{22}$ 个氯离子。一个氯离子的直径约 $3.6\times 10^{-10}$ 米。如果把这一勺盐里的氯离子排成一行，总长度大约是多少米？和地球到月球的距离（约 $3.8\times 10^{8}$ 米）比一比。
\item 用能量账的语言重写给氩气球“验明正身”的理由：为什么节日气球充氦气比充氢气安全？（提示：分别算“从它身上抢电子”和“往它身上塞电子”两笔账，再想想氢气的账长什么样。）
\item 画一张卤钨循环的物质流图：用方框表示“灯丝上的钨”“蒸发的钨”“碘化钨”，用箭头表示它们之间的转变，并在每条箭头旁标出该处的温度是高还是低。图旁注明：方框和箭头是人为的表示法。
\end{sikaoti}
